\statement{C.38}{Proved}{The paired involution has exactly the peripheral coinvariant kernel}
Retain \(\alpha,u,v\), and \(\tau=\alpha v\) from C.37, and put
\(z=(xy)^{-1}\). Let
\[
 N_\tau=\overline{\langle\!\langle
          \tau(g)g^{-1}:g\in F\rangle\!\rangle}_F,
 \qquad
 N_z=\overline{\langle\!\langle z\rangle\!\rangle}_F .
\]
Then
\[
             N_\tau=N_z,\qquad F/N_\tau\simeq\widehat{\mathbb Z}.
 \tag{C.38.1}
\]
Moreover there are a local involutory section
\(r\in F\rtimes_\tau C_2\) that inverts \(z\), and \(g\in F\),
such that
\[
                              z=[g,r].
 \tag{C.38.2}
\]
Equivalently \(z\) is \(F\)-conjugate to a single twisted
commutator \(h\tau(h)^{-1}\). These equalities do not assert
that \(g,r\) generate the semidirect product, or identify that
group with a free profinite product.

\begin{proof}
The involution \(\tau\) preserves the conjugacy class of
\(z^{-1}\), so it preserves \(N_z\). The quotient by \(N_z\)
has the presentation \(y=x^{-1}\), and hence is the full
procyclic group \(\widehat{\mathbb Z}\). The action of
\(\tau\) on \(F^{\rm ab}\) is minus interchange, which becomes
identity after imposing \(y=x^{-1}\). Thus \(\tau\) acts
identically on \(F/N_z\), proving \(N_\tau\subseteq N_z\).

Conversely take any finite quotient of \(F/N_\tau\).
The induced action of \(\tau\) is identity by definition.
The one-section consequence (C.37.4) makes the image of
\(z\) trivial in every such quotient. A profinite group is
separated by its finite quotients, so \(z\in N_\tau\).
Normality gives \(N_z\subseteq N_\tau\), proving (C.38.1).

Choose the local section \(r\) inverting \(z\) exactly as
in the proof of (C.37.4). Both \(r\) and \(zr\) are
involutory sections. There is a single \(F\)-conjugacy class
of sections, so \(zr=grg^{-1}\) for some \(g\in F\).
Since \(r^2=1\), this gives (C.38.2). Also \(r\) is
\(F\)-conjugate to the specified section \(\tau\).
Conjugating (C.38.2) accordingly gives
\([h,\tau]=h\tau(h)^{-1}\), which proves the last assertion.
No generation conclusion follows from this conjugacy alone.
\end{proof}
